\subsection*{Independent second reading of the normal-subgroup criterion}

All 425 lines of
\texttt{arithmetic\_notes/B8\_normal\_extensions.tex}
were independently read on 18 November 2026.
The outcome for B.META.17 and B.META.18, assembled as
B.142 and B.143, is \textbf{PASS}; no correction was
requested.  The exact group-theoretic source passages
were independently reopened: Kochloukova--Zalesskii,
Theorem~5.1, Case~I, arXiv version printed p.~7, and
Iwasawa, Theorem~2(ii) and its preceding paragraph,
printed p.~463.  The AMS retrieval of the latter
returned 403; the actual journal scan at the recorded
SciSpace PDF URL supplied the opened primary text.
This is a same-model-family reading by GPT-6 Astra Pro
in Pro mode, with no separately exposed effort setting.

\medskip\noindent\textbf{B.142: precisely the arithmetic quotient that is needed.}
The hypotheses give $\mu_\infty\subset A\subset K$.
For finite $E/K$, adjoining one finite Galois $D_1/D$
containing its generators and the element under study
leaves the local Kummer argument in a single finite
local field.  The groups $H_1$, $\Gamma$ and $C$ in
(META.NG3) have the asserted normality, and $C$ is
abelian with trivial Tate coefficient action.

For each $g\in B$, conjugation is trivial on $C$
because $[B,B]\subset H$, and it is trivial on
$\Gamma/C$ because $B$ is normal in inertia.
Its action on the coefficients is also trivial.
Inflation--restriction therefore makes $(g-1)^2$
zero on the relevant degree-one cohomology.  The
finite orbit of the specific algebraic element makes
its Kummer vector fixed by a positive power of $g$;
over $\mathbb Q_p$ this forces that vector to be
fixed by $g$.  The local restriction injectivity
and the finite-local-field Kummer kernel then put
all ratios $g(a)/a$ in the one finite group of
roots of unity.  One exponent works for all $g$.

Only this implication is proved and used.  In
particular the proof never assumes that every
radical of $A$ lies in $K$ or has zero Kummer
class over it.  Strict containment $H\subsetneq B$
provides $b\in K\setminus A$, and the elementary
two-element identity for $b,1-b$ produces a
nonzero stable class from the one-sided implication
alone.  The two-grid argument requires the same
one-sided implication, not equality of the kernel
with a radical subgroup.

The remaining norm-proof inputs are also checked,
rather than inferred merely from the kernel.
The normalized arithmetic constants come from
actual values in $D_1$; normality of $H$ gives the
required local Galois invariance.  The chosen
surjective integral Kummer class yields arbitrary
cyclic $p$-power local extensions, still with the
normal kernel fixed.  B.139 supplies algebraic
local bases and dense coordinates.  The quotient
inflation argument gives a scalar local logarithmic
error.  Nonzero stable Kummer classes ensure the
non-separably-closed hypothesis required by the
opened etale-open/valuation comparison.

The subsequent use of B.136 is the already audited
norm reconstruction: analytic reduced ratios,
root types and their divisor cocycle, exclusion of
constant roots, positivity by specialization,
finite-grid rationality, global removal of the
common factor, graph degree one, and continuity
at exceptional points.  These arguments apply
over $\Omega$ exactly as verified in B.140.
Every construction starts with an arbitrary finite
extension of $K$, so the final reconstruction
criterion's finite-extension quantifier is met.

\medskip\noindent\textbf{B.143: no nontrivial abelian normal local subgroup.}
If a nontrivial abelian normal subgroup $N$ met
wild inertia nontrivially, take an element of
that intersection together with two noncommuting
wild elements.  Their closed generated subgroup
is a finitely generated nonabelian free pro-$p$
group.  Its intersection with $N$ is a nontrivial
abelian normal subgroup; it lies in the procyclic
centralizer of the chosen element and is therefore
procyclic.  The exact normalizer statement in
Kochloukova--Zalesskii, Case~I, forces the entire
generated group to be procyclic, a contradiction.

Thus the intersection is trivial.  Normality makes
$[N,P]$ lie in that intersection, so $N$ centralizes
wild inertia.  The local centralizer theorem
proved in B.META.1 then makes $N$ trivial.  This
argument is valid at every prime and for every
finite local extension.  Iwasawa's precise input
is the countably infinite free pro-$p$ wild group;
the finite-generator reduction respects the scope
of the cited normalizer theorem.

\medskip\noindent\textbf{The subgroup $H$ is separately nonzero and proper.}
For arbitrary nontrivial normal $N$, put $B=N'$.
The previous argument proves $B\ne1$, and $B$
is still normal in the whole local group.  The
unramified quotient and the cyclotomic image are
abelian, so $B\subset I\cap\ker\chi_p$.
There is no claim that this $B$ is pro-$p$.

Finite inertia quotients are solvable, and hence
the intersection of its finite derived series
is trivial.  The least $r\ge1$ such that
$B\not\subset I^{(r)}$ therefore exists and gives
$B\subset I^{(r-1)}$.  The intersection
$H=B\cap I^{(r)}$ is consequently proper in $B$
and contains $B'$.  The separate nonabelian-normal
subgroup result, applied to $B$, gives $B'\ne1$.
Thus $H\ne1$.  This is the essential second
argument; properness alone would not suffice.
The intersection is closed and normal in the
whole local Galois group, and $B/H$ is nontrivial
abelian.  All hypotheses of B.142 are now met.

Since $H\subset N$, the resulting reconstructed
field contains the original fixed field as an
algebraic subfield.  The explicit descent argument
of B.140 applies both directly and after every
finite compositum.  The asserted finite-extension
and algebraic-subfield conclusions follow.

\medskip\noindent\textbf{The core condition is used literally.}
In the more general criterion, the indicated
intersection of conjugates is assumed nontrivial.
It is a closed normal subgroup of the actual
open local group $G_{M_1}$ and is contained in
$G_F$.  B.139 realizes $M_1$ as the completion
of a finite extension of the original finitely
generated field, so the preceding result applies
with this new base.  Its fixed field contains
$F$, and algebraic descent gives the conclusion.
No normal closure is substituted and no
nontriviality is inferred from the original
subgroup merely being nontrivial.

Finally, the reduction of a henselian subfield
already embedded in $\overline{\mathbb Q}_p$
uses a finitely generated subfield containing a
transcendence basis and the universal property
of its chosen henselization.  It does not claim
that an arbitrary mixed-characteristic henselian
field embeds in this algebraic closure.
